% Videos: network security model, symmetric cipher model, types of crypto, security of schemes
\section{Model for Network Security}

Sender $S$ wants to send a message to receiver $R$ over an insecure channel (the Internet). An \textbf{opponent} can sniff the channel.

\begin{center}
\begin{tikzpicture}[node distance=6mm]
  \node[box,fill=cyan!10] (m1) {Message\\(S)};
  \node[circ,right=of m1] (t1) {T};
  \node[box,below=5mm of t1,fill=magenta!8] (k1) {Secret info};
  \node[box,right=of t1,fill=yellow!20] (c1) {Secure\\message};
  \node[box,right=8mm of c1,minimum width=26mm,fill=gray!15] (ch) {Information\\channel};
  \node[box,right=8mm of ch,fill=yellow!20] (c2) {Secure\\message};
  \node[circ,right=of c2] (t2) {T};
  \node[box,below=5mm of t2,fill=magenta!8] (k2) {Secret info};
  \node[box,right=of t2,fill=cyan!10] (m2) {Message\\(R)};
  \node[box,below=6mm of ch,fill=red!10] (op) {Opponent};
  \node[box,above=7mm of ch,fill=green!12] (ttp) {Trusted Third Party};
  \draw[arr] (m1)--(t1); \draw[arr] (k1)--(t1); \draw[arr] (t1)--(c1);
  \draw[arr] (c1)--(ch); \draw[arr] (ch)--(c2);
  \draw[arr] (c2)--(t2); \draw[arr] (k2)--(t2); \draw[arr] (t2)--(m2);
  \draw[darr] (op)--(ch);
  \draw[darr,orange!80!black] (ttp)--(ch);
  \draw[arr,orange!80!black,dashed] (ttp.west) -| (k1.north west);
  \draw[arr,orange!80!black,dashed] (ttp.east) -| (k2.north east);
\end{tikzpicture}
\end{center}

\begin{itemize}
  \item \textbf{T (security-related transformation)} takes two inputs: the message and \textbf{secret info} (a key). Output: a secure message (ciphertext). $R$ applies the reverse $T$ with its secret info.
  \item It is called $T$, not $E$, because encryption is only one option. \textbf{Steganography} (hiding data inside an image, audio, etc.) is another.
  \item The opponent can intercept the secure message but cannot read it.
  \item Security rests on the \textbf{key} (its secrecy and bit length) and algorithm strength, not on hiding the algorithm. Small key space $\Rightarrow$ brute force works.
  \item \textbf{Trusted Third Party (TTP)}: distributes keys, issues certificates, settles disputes. Example: a Certificate Authority signing a bank's SSL certificate. (Details come later in the course.)
\end{itemize}

\subsection{Four tasks in designing a security service}
\begin{keybox}
\begin{enumerate}[leftmargin=1.5em]
  \item \textbf{Design an algorithm} for $T$ (encrypt at $S$, decrypt at $R$) that the opponent cannot defeat.
  \item \textbf{Generate the secret information} (large keys, so brute force is infeasible).
  \item \textbf{Develop methods to distribute and share} the secret info (get keys to $S$ and $R$ safely).
  \item \textbf{Specify a protocol}: the rules both parties follow to use the algorithm and keys (e.g., SSL/TLS, IPsec).
\end{enumerate}
\end{keybox}
Memory aid: \emph{Algorithm $\to$ Key $\to$ Distribute $\to$ Protocol}.

\section{Network Access Security Model}

Protects a \emph{system} from unwanted access, rather than protecting a message in transit.

\begin{center}
\begin{tikzpicture}[node distance=8mm]
  \node[box,fill=red!10,text width=30mm] (op) {\textbf{Opponent}\\Human (attacker)\\Software (virus, worm)};
  \node[box,right=of op,fill=gray!15,minimum width=24mm] (ac) {Access channel};
  \node[box,right=of ac,fill=green!12,text width=22mm] (gk) {\textbf{Gatekeeper function}\\(firewall)};
  \node[box,right=12mm of gk,fill=cyan!8,text width=42mm,align=left] (is) {\textbf{Information System}\\
    \footnotesize Computing resources (CPU, memory, I/O)\\ \footnotesize Data \quad Processes \quad Software\\
    \footnotesize\textbf{Internal security controls} (antivirus)};
  \draw[arr] (op)--(ac); \draw[arr] (ac)--(gk); \draw[arr] (gk)--node[above,font=\scriptsize]{allowed} (is);
\end{tikzpicture}
\end{center}

\begin{itemize}
  \item \textbf{Opponent}: a human hacker or malicious software.
  \item \textbf{Gatekeeper}: first line of defence, like a guard at a gate. Filters traffic and blocks intruders (firewall, login checks).
  \item \textbf{Internal security controls}: second line, inside the system, catching what slips past the gate (antivirus, OS permissions).
  \item Insiders (e.g., employees) can also attack, but this model focuses on outside threats.
\end{itemize}

\section{Cryptography}

\begin{keybox}[title=Definition]
The art or science of transforming an \textbf{intelligible} message (plaintext) into an \textbf{unintelligible} one (ciphertext), and then \textbf{retransforming} it back to the original form.
\end{keybox}

\subsection{Symmetric cipher model}
\begin{center}
\begin{tikzpicture}[node distance=7mm]
  \node[box,fill=cyan!10] (p1) {Plaintext $P$};
  \node[box,right=of p1,fill=green!15] (e) {Encryption\\algorithm $E$};
  \node[box,right=12mm of e,fill=yellow!20] (c) {Ciphertext $C$};
  \node[box,right=16mm of c,fill=red!12] (d) {Decryption\\algorithm $D$};
  \node[box,right=of d,fill=cyan!10] (p2) {Plaintext $P$};
  \node[box,above=5mm of e,fill=yellow!40] (k1) {Key $K$};
  \node[box,above=5mm of d,fill=yellow!40] (k2) {Key $K$};
  \draw[arr] (p1)--(e); \draw[arr] (k1)--(e); \draw[arr] (e)--(c);
  \draw[arr] (c)--node[above,font=\scriptsize]{transmitted}(d); \draw[arr] (k2)--(d); \draw[arr] (d)--(p2);
  \draw[darr,bend left=15,orange!80!black] (k1.north) to node[above,font=\scriptsize]{same key} (k2.north);
\end{tikzpicture}
\end{center}
\[
C = E(K, P) \qquad\qquad P = D(K, C) = D\big(K, E(K,P)\big)
\]
\begin{itemize}
  \item Encryption needs \textbf{two} inputs ($P$ and $K$); decryption needs $C$ and $K$.
  \item $E$ and $D$ are \textbf{public}. Only $K$ is secret (Kerckhoffs's principle).
  \item If $K$ leaks, \emph{all} past and future messages between $S$ and $R$ can be read.
  \item Toy example: $K = 100$ at both ends. Real keys are huge (128, 256, 2048 bits), so guessing is infeasible.
\end{itemize}

\subsection{Types of cryptography}
\begin{center}
\small
\begin{tabularx}{\textwidth}{@{}l X X@{}}
\toprule
 & \textbf{Symmetric} (private / secret key) & \textbf{Asymmetric} (public key) \\
\midrule
Keys & One shared key & Key pair: public + private \\
Use & Same key encrypts and decrypts & One key encrypts, the other (mathematically related) decrypts \\
Secrecy & Key known only to $S$ and $R$ & Public key published; private key kept by owner \\
Toy example & Encrypt 100, decrypt 100 & Encrypt 100, decrypt 500 \\
\bottomrule
\end{tabularx}
\end{center}
Asymmetric works because the two keys satisfy a mathematical relationship, and the private key cannot practically be computed from the public one. (Which key does what is covered in a later lecture.)

\section{Security of Encryption Schemes}\label{sec:schemes}

\begin{keybox}[title=Unconditionally secure]
The ciphertext does \textbf{not contain enough information to uniquely determine the plaintext}, no matter how much ciphertext is available. Holds even against infinite computing power. Example: one-time pad.
\end{keybox}

\begin{keybox}[title=Computationally secure (either condition)]
\begin{enumerate}[leftmargin=1.5em]
  \item \textbf{Cost} of breaking the cipher $>$ \textbf{value} of the encrypted information.
  \item \textbf{Time} to break the cipher $>$ \textbf{useful lifetime} of the information.
\end{enumerate}
\end{keybox}
\begin{itemize}
  \item Cost: spending \$1{,}000{,}000 to steal a \$1{,}000 message is pointless.
  \item Time: a quote valid for 10 minutes that takes 50 years to crack is useless when cracked.
  \item Real ciphers (AES, RSA) are computationally secure; unconditional security needs a key as long as the message, which is impractical.
\end{itemize}

\begin{examtip}
\begin{itemize}
  \item Draw the network security model with all labels: $S$, $R$, $T$, secret info, secure message, opponent, TTP.
  \item Know the four tasks in order.
  \item ``Why $T$ and not $E$?'' $\Rightarrow$ steganography is also a transformation.
  \item Symmetric = same key; asymmetric = different but related keys.
  \item Quote both conditions of computational security word for word.
\end{itemize}
\end{examtip}
